% Optional strengthening of the ultrapower independence construction.
\subsection{A stronger initial structure with cancellative addition}

The initial structure can already satisfy additive cancellation and the
identity law for successor iteration. Take
\[
 M_0=\{(0,n):n\in\mathbb N\}
     \cup\{(k,z):k\in\mathbb N_{>0},\ z\in\mathbb Z\},
\]
with componentwise addition, $0=(0,0)$, $1=(0,1)$, and
$S(k,z)=(k,z+1)$. It is a submonoid of the additive group $\mathbb Z^2$,
so addition is associative, commutative, and cancellative, with identity
$0$. It is conical: $a+b=0$ implies $a=b=0$. Successor is injective,
omits zero, has no finite cycle, and every nonzero element has a
predecessor in the domain.

Preinstall two iterator symbols:
\[
 J(n,a)=
 \begin{cases}
 a&\text{if the first coordinate of }n\text{ is }0,\\
 1&\text{otherwise},
 \end{cases}
 \qquad K(n,a)=a+n.
\]
They represent the updates $u\mapsto u$ and $u\mapsto Su$, respectively.
Both satisfy their zero and successor computations. The added identity
law $K(n,0)=n$ holds, as do the usual additive zero and successor laws.
Their cross-extensionality condition has a false antecedent, since
identity and successor are distinct functions. Extensionality within
each description is immediate. Thus this is a valid starting point for
the same extension tower. All these first-order laws survive that tower.

At $c=(1,0)$, however, $J(c,0)=1$. Hence the final countermodel can have
all of the listed additive laws, including cancellation, and the
successor identity orbit, while still refuting literal zero-equation
induction.

This does not furnish a countermodel to full successor uniqueness.
Indeed, already in the initial structure the term
\[
 f(n)=n+J(n,0)
\]
satisfies $f(0)=0$ and $f(Sn)=Sf(n)$, whereas
$f(c)=Sc\ne c=K(c,0)$. Its addition is itself the successor-iterator
operation $K$, so this failure uses scalar terms. Full successor
uniqueness therefore fails in this structure, as it must.
